\textbf{Время:} 01:15

\textbf{Тезис:} пространственная разность q описывает относительное движение, а полусумма K — общее.

\textbf{Визуал:} два вектора импульса; q соединяет их концы, K направлен между ними.

\textbf{Текст:} Пока достаточно обычных трёхмерных векторов. Вычтем второй импульс из первого и назовём разность q. Близкие векторы дают малую разность. Не путайте близость импульсов с близостью направлений: одинакового угла недостаточно, если модули сильно отличаются.

Полусумму назовём K. Она описывает общее движение пары. Поперечный импульс k-T — модуль поперечной части этой векторной полусуммы. Это не среднее двух модулей: противоположные одинаковые векторы дают ноль.

Во всём рассказе q означает полную разность. Релятивистскую запись с энергией добавим тогда, когда она понадобится для фазы. Сейчас уже можно разбивать пары по относительному и среднему импульсам.

\textbf{Переход:} Теперь у нас есть переменные для пары. Объясним квантовую причину особой корреляции при малом q.
